{\centering
	\chapter{Infinite Go}\label{infgo}}
\thispagestyle{empty}
\noindent
{\large\textbf{This chapter it's a draft, at least for now. I'll include my thought process and most of what I'm able to understand/formalize about it.}}
\vspace{20pt}

To begin with, let's see \textit{where} we are playing now. Our infinite board should be able to have also infinite dimensions and every one of those could be or not infinite as well.

Let's begin with the simplest case (I hope), let's play on a two dimensional board that has one infinite dimension.

Let $G$ be a go board such that,
$$G=\mathbb{N}\times N.$$
where $N$ is a subset of $\mathbb{N}$.
As the number of players does not matter at all, let's use only two players and, the canonical adjacency.
Our game will be $\Gamma=(G,E,A_{ort},\Delta_{Go_{\infty}})$.
At least for now, we are not sure how $\Delta_{Go_{\infty}}$ works but, we will get there.

Our board could (and for us, will) look like this, Figure~\ref{first_inf_board}:
\begin{figure}[H]
	%\begin{figure}[H]
	\centering
	\begin{tikzpicture}[scale=0.485] \tikzstyle{every node}=[font=\small]
		% Draw the background
		\fill[goban] (-7,-0.5) rectangle (7,6.5);
		% Draw the grid lines
		\foreach \x in {6,5,...,0}{
			\draw[black] (-5,\x) -- (5,\x); 
			% draw horizontal lines
		}
		\foreach \x in {4,...,-4} {
			\draw[black] (\x,0) -- (\x,6); 
			% draw vertical lines
		}
		\foreach \x in {6,0}{
		\draw[very thick] (-5,\x) rectangle (5,\x);
		}% draw upper and lower lines.
		%draw the triple dot line for each infinite side.
		\foreach \x in {6,5,...,0} {
			\node at (-6, \x) {\dots};
			\node at (6, \x) {\dots};
		}
		
	
	\end{tikzpicture}
	\caption{A $7\times \infty$ board.}
	\label{first_inf_board}
\end{figure}

The very first problem on the infinite board it's how to determine a capture. Any other prior definition should work just fine.

Let's say we are in the next position Figure~\ref{first_inf_cap}
\begin{figure}[H]
	%\begin{figure}[H]
	\centering
	\begin{tikzpicture}[scale=0.485] \tikzstyle{every node}=[font=\small]
		% Draw the background
		\fill[goban] (-7,-0.5) rectangle (7,6.5);
		% Draw the grid lines
		\foreach \x in {6,5,...,0}{
			\draw[black] (-5,\x) -- (5,\x); 
			% draw horizontal lines
		}
		\foreach \x in {4,...,-4} {
			\draw[black] (\x,0) -- (\x,6); 
			% draw vertical lines
		}
		\foreach \x in {6,0}{
			\draw[very thick] (-5,\x) rectangle (5,\x);
		}% draw upper and lower lines.
		%draw the triple dot line for each infinite side.
		\foreach \x in {6,5,...,0} {
			\node at (-6, \x) {\dots};
			\node at (6, \x) {\dots};
		}
		
		%draw black stones.
		\foreach \x in {5,...,-5} {
			\filldraw[black] (\x,3) circle (12pt);
		}
		\begin{scope}[xshift=15cm]
			% Draw the background
			\fill[goban] (-7,-0.5) rectangle (7,6.5);
			% Draw the grid lines
			\foreach \x in {6,5,...,0}{
				\draw[black] (-5,\x) -- (5,\x); 
				% draw horizontal lines
			}
			\foreach \x in {4,...,-4} {
				\draw[black] (\x,0) -- (\x,6); 
				% draw vertical lines
			}
			\foreach \x in {6,0}{
				\draw[very thick] (-5,\x) rectangle (5,\x);
			}% draw upper and lower lines.
			%draw the triple dot line for each infinite side.
			\foreach \x in {6,5,...,0} {
				\node at (-6, \x) {\dots};
				\node at (6, \x) {\dots};
			}
			
			%draw black stones.
			\foreach \x in {5,...,-5} {
				\filldraw[draw=black,fill=white] (\x,4) circle (12pt);
			}
			%draw black stones.
			\foreach \x in {5,...,-5} {
				\filldraw[black] (\x,3) circle (12pt);
			}
			%draw black stones.
			\foreach \x in {5,...,-5} {
				\filldraw[draw=black,fill=white] (\x,2) circle (12pt);
				
			}
		\end{scope}
		
	\end{tikzpicture}
	\caption{An infinite board with an infinite black group being capture.}
	\label{first_inf_cap}
\end{figure}
As this group is infinite (numerable), will touch the very end of the board so we can capture it by surronding it at the top and bottom.

We can be sure that this is capture because if the black group it's $I$, then $K(I)\subset S$ where $S$ is the set of the two white groups. That is, $L(I)$ should be $0$.
Now let's take a look at another example, an infinite group but only on one side.
\begin{figure}[H]
	%\begin{figure}[H]
	\centering
	\begin{tikzpicture}[scale=0.485] \tikzstyle{every node}=[font=\small]
		% Draw the background
		\fill[goban] (-7,-0.5) rectangle (7,6.5);
		% Draw the grid lines
		\foreach \x in {6,5,...,0}{
			\draw[black] (-5,\x) -- (5,\x); 
			% draw horizontal lines
		}
		\foreach \x in {4,...,-4} {
			\draw[black] (\x,0) -- (\x,6); 
			% draw vertical lines
		}
		\foreach \x in {6,0}{
			\draw[very thick] (-5,\x) rectangle (5,\x);
		}% draw upper and lower lines.
		%draw the triple dot line for each infinite side.
		\foreach \x in {6,5,...,0} {
			\node at (-6, \x) {\dots};
			\node at (6, \x) {\dots};
		}
		
		%draw black stones.
		\foreach \x in {5,...,0} {
			\filldraw[black] (\x,3) circle (12pt);
		}
		\begin{scope}[xshift=15cm]
			% Draw the background
			\fill[goban] (-7,-0.5) rectangle (7,6.5);
			% Draw the grid lines
			\foreach \x in {6,5,...,0}{
				\draw[black] (-5,\x) -- (5,\x); 
				% draw horizontal lines
			}
			\foreach \x in {4,...,-4} {
				\draw[black] (\x,0) -- (\x,6); 
				% draw vertical lines
			}
			\foreach \x in {6,0}{
				\draw[very thick] (-5,\x) rectangle (5,\x);
			}% draw upper and lower lines.
			%draw the triple dot line for each infinite side.
			\foreach \x in {6,5,...,0} {
				\node at (-6, \x) {\dots};
				\node at (6, \x) {\dots};
			}
			
			%draw white stones.
			\foreach \x in {5,...,0} {
				\filldraw[draw=black,fill=white] (\x,4) circle (12pt);
			}
			%draw black stones.
			\foreach \x in {5,...,0} {
				\filldraw[black] (\x,3) circle (12pt);
			}
			%draw white stones.
			\foreach \x in {5,...,0} {
				\filldraw[draw=black,fill=white] (\x,2) circle (12pt);
				
			}
			\filldraw[draw=black,fill=white] (-1,3) circle (12pt);
		\end{scope}
		
	\end{tikzpicture}
	\caption{An infinite, on one side, black group being capture.}
	\label{first_1sideinf_cap}
\end{figure}
In this case, not using the whole board, even though it's infinite, its better. As it provides one more liberty (or does not loses it).

Let's see the following example, we might into the nature of inifinite Go.

\begin{figure}[H]
	%\begin{figure}[H]
	\centering
	\begin{tikzpicture}[scale=0.485] \tikzstyle{every node}=[font=\small]
		% Draw the background
		\fill[goban] (-7,-0.5) rectangle (7,6.5);
		% Draw the grid lines
		\foreach \x in {6,5,...,0}{
			\draw[black] (-5,\x) -- (5,\x); 
			% draw horizontal lines
		}
		\foreach \x in {4,...,-4} {
			\draw[black] (\x,0) -- (\x,6); 
			% draw vertical lines
		}
		\foreach \x in {6,0}{
			\draw[very thick] (-5,\x) rectangle (5,\x);
		}% draw upper and lower lines.
		%draw the triple dot line for each infinite side.
		\foreach \x in {6,5,...,0} {
			\node at (-6, \x) {\dots};
			\node at (6, \x) {\dots};
		}
		
		%draw black stones.
		\foreach \x in {1,...,-1} {
			\filldraw[black] (\x,3) circle (12pt);
		}
		\foreach \x in {2,4} {
			\filldraw[black] (0,\x) circle (12pt);
		}
		\filldraw[draw=black,fill=white] (0,1) circle (12pt);
		\filldraw[draw=black,fill=white] (0,5) circle (12pt);
		\filldraw[draw=black,fill=white] (-1,4) circle (12pt);
		\filldraw[draw=black,fill=white] (-1,2) circle (12pt);
		\filldraw[draw=black,fill=white] (1,4) circle (12pt);
		\filldraw[draw=black,fill=white] (1,2) circle (12pt);
		\filldraw[draw=black,fill=white] (2,3) circle (12pt);
		\filldraw[draw=black,fill=white] (-2,3) circle (12pt);

		
		\begin{scope}[xshift=15cm]
			% Draw the background
			\fill[goban] (-7,-0.5) rectangle (7,6.5);
			% Draw the grid lines
			\foreach \x in {6,5,...,0}{
				\draw[black] (-5,\x) -- (5,\x); 
				% draw horizontal lines
			}
			\foreach \x in {4,...,-4} {
				\draw[black] (\x,0) -- (\x,6); 
				% draw vertical lines
			}
			\foreach \x in {6,0}{
				\draw[very thick] (-5,\x) rectangle (5,\x);
			}% draw upper and lower lines.
			%draw the triple dot line for each infinite side.
			\foreach \x in {6,5,...,0} {
				\node at (-6, \x) {\dots};
				\node at (6, \x) {\dots};
			}
			
			
			%draw white stones.
			\foreach \x in {5,...,-5} {
				\filldraw[draw=black,fill=white] (\x,4) circle (12pt);
			}
			%draw white stones.
			\foreach \x in {5,...,-5} {
				\filldraw[draw=black,fill=white] (\x,2) circle (12pt);
				
			}
			%draw white stones.
			\foreach \x in {0,1,2,4,5,6} {
				\filldraw[draw=black,fill=white] (-1,\x) circle (12pt);
				
			}
			\foreach \x in {0,1,2,4,5,6} {
				\filldraw[draw=black,fill=white] (1,\x) circle (12pt);
				
			}
			
			%draw black stones.
			\foreach \x in {5,...,-5} {
				\filldraw[black] (\x,3) circle (12pt);
			}
			\foreach \x in {0,...,6} {
				\filldraw[black] (0,\x) circle (12pt);
			}
		\end{scope}
		
	\end{tikzpicture}
	\caption{Two captures.}
	\label{finite_infinite_capture}
\end{figure}
It seems that finite groups are stronger than infinite ones as, again, they don't lose the liberty due to being no infinite.

\section{SuperKo at infinity}

We cannot use Zobrits method for the infinite Go as we will need infinite many bit's. 
We need a better method, let's try to make it.

If we remember the nature of the SuperKo rule, it's says that we cannot replicate a prior board via a stone play. Again, we are using this version of SuperKo.

Because of the rule, every game has to be a line, that is, as we cannot go back into the historial (at least it's a pass but will not take in consideration those). From the empty board to the very end, we have a straight forward path. Not even that but, because of the rule, the whole game of Go its a \textit{tree}. That is, starting from the empty board, there are 362 possbile plays (361 stone placements and a pass) and we cannot go back.

Let's define a way to measure how close are two given boards (finite).
\begin{definition}[Color difference function]
	Given two tables $T_1$ and $T_2$ of the same Go game, the difference in a point $t=(t_1,...t_m)$, we define the color difference as the following,
	\[ 
	\delta(T_1,T_2)_t=
	\begin{cases}
		1 \text{if }C(t)_{T_1}=C(t)_{T_2}, \\
		0 \text{if } not. \\
	\end{cases}
	\]
\end{definition}
Now, let's apply this to every point,
\begin{definition}[Distance between Tables]
	Given two tables $T_1$ and $T_2$ of the same Go game, the difference between them will be:
	$$\sum_{t\in G} \delta(T_1,T_2)_{t}$$
\end{definition}
I have no idea how to use this, yet.